\documentclass[11pt,a4paper]{article}
\usepackage[utf8]{inputenc}
\usepackage[T1]{fontenc}
\usepackage[german]{babel}
\usepackage{amsmath,amssymb,amsfonts}
\usepackage{geometry}
\geometry{margin=2.5cm}
\usepackage{microtype}
\usepackage{color}

\title{\textbf{Synthese des EABC-Modells}\\ \large Von quaternionischen Bulks zu holografischer Vakuumimpedanz, $E_8$-Gittertopologie und dem kosmischen Druckabfall}
\author{Thomas Hoffbauer \thanks{Forschungsbericht im Rahmen des EABC-Modells (Bamberg Model)}}
\date{September 2026}

\begin{document}

\maketitle

\begin{abstract}
Traditionelle Ansätze der theoretischen Physik leiden unter dimensionsbehafteten Naturkonstanten und unkontrollierten ultravioletten Divergenzen. Das EABC-Modell (Euler-Adic-Borel-Cantor) löst dieses Dilemma durch eine konsequente topologische Geometrisierung. Durch die fundamentalen Setzungen $\hbar c = \pi$, die Gravitationskonstante an der Horwitz-Kugel $G = 4$ und die Einbettung in eine Typ-$\text{II}_1$-von-Neumann-Algebra entpuppen sich physikalische Konstanten als rein algebraische Invarianten. Der Casimir-Effekt, die Feinstrukturkonstante $\alpha$ und die Hubble-Spannung werden dabei als holografische Manifestationen des $E_8$-Wurzelsystems und der Gabriel-Hörner identifiziert.
\end{abstract}

\section{1. Das arithmetische Fundament ($\hbar c = \pi$ und $G = 4$)}
Im Zentrum des EABC-Modells steht der Verzicht auf empirische Fit-Parameter. Die Raumzeit wird über vier orthogonale Tori aufgespannt, die auf den Hamilton-Hurwitz-Quaternionen $\mathbb{H} = \{1, i, j, k\}$ basieren. 
\begin{itemize}
    \item \textbf{Die Wirkungs-Kreis-Identität:} Die Kopplung von Wirkungsquantum und Lichtgeschwindigkeit wird auf die grundlegende Kreiszahl geeicht:
    \begin{equation}
        \hbar c = \pi
    \end{equation}
    \item \textbf{Die Gravitations-Basis:} An der Oberfläche der Horwitz-Kugel ist die Gravitationskonstante starr durch die vier quaternionischen Dimensionen gegeben:
    \begin{equation}
        G = 4
    \end{equation}
\end{itemize}
Aus dieser Kombination emergieren die Planck-Einheiten (wie die Planck-Masse $m_P = \frac{\sqrt{\pi}}{2}$) nicht als absolute Schranken, sondern als reine geometrische Eigenwerte.

\section{2. Bulk-Rand-Dualität und die Geometrische Impedanz}
Das Universum gliedert sich topologisch in ein endliches inneres Volumen (den Bulk mit den Primzahlquadrupeln $\mathcal{P}_4$) und einen asymptotischen Rand (die Gabriel-Hörner bzw. Aharonov-Bohm-Solenoide).
\begin{itemize}
    \item \textbf{Die Geometrische Impedanz:} Der Übergang zwischen Bulk und Rand wird durch die Impedanz des Gabriel-Horns reguliert. Sie fungiert als natürlicher Dämpfungsfaktor, der das Oberflächen-Paradoxon unendlicher Hohlkörper auflöst.
    \item \textbf{Typ-$\text{II}_1$-Spurregulierung:} Die Renormierung erfolgt über die Fuglede-Kadison-Determinante einer von-Neumann-Algebra vom Typ $\text{II}_1$, wodurch Divergenzen in der QED und Yang-Mills-Theorie (Mass Gap) arithmetisch gebändigt werden.
\end{itemize}

\section{3. Holografischer Casimir-Druck und das $E_8$-Gitter}
Die Anwendung auf den Vakuumdruck (Casimir-Effekt) offenbart die tiefste Symmetrie des Modells. Unter Berücksichtigung von $\hbar c = \pi$ reduziert sich die klassische Druckformel auf:
\begin{equation}
    P_{\text{Casimir}} = -\frac{\pi^3}{240 d^4}
\end{equation}
\begin{enumerate}
    \item \textbf{Der Zähler ($\pi^3$):} Repräsentiert das dreidimensionale holografische Volumen des inneren Torus-Bulks.
    \item \textbf{Der Nenner ($240$):} Korrespondiert exakt mit den \textbf{240 Wurzelvektoren des $E_8$-Wurzelsystems} in 8 Dimensionen.
    \item \textbf{Der Raum-Term ($d^4$):} Verankert den Druck in den vier Dimensionen der Quaternionen-Metrik.
\end{enumerate}

\section{4. Die Hubble-Spannung als makroskopischer Druckabfall}
Der mikroskopische Casimir-Druck ist identisch mit dem globalen Verhalten des Kosmos. Die Indexhöhe $n$ der Primzahlquadrupel fungiert als Renormierungsgruppen-Skala (Temperatur- und Druckskala). 
\begin{itemize}
    \item Die \textbf{Hubble-Spannung ($\Delta H_0$)} ist folglich kein Messfehler, sondern der messbare \textbf{Druckabfall} beim Übergang vom hochenergetischen Bulk-Gitter des frühen Universums zum asymptotischen Rand des späten Kosmos.
    \item Die verbleibende kosmische Hintergrundspannung (Dunkle Energie) wird durch die invariante Rand-Spur der Gabriel-Hörner garantiert.
\end{itemize}

\section{5. Fazit}
Das EABC-Modell beweist, dass Materie, Eichfelder und Kosmologie Projektionen einer reinen Zahlentheorie sind. Primzahlen, Quaternionen und $E_8$-Gitter bilden das Betriebssystem einer raumzeitlichen Realität, deren Konstanten universell aus der Geometrie des Kreises ($\pi$) abgelesen werden können.

\end{document}